\section{Especificación del Transformer mínimo}
\label{sec:transformer_minimo}

Las secciones anteriores establecieron los fundamentos teóricos del
Transformer: el mecanismo de atención, la arquitectura completa y sus
variantes. Esta sección define formalmente el modelo concreto que se
implementa (ver el repositorio), especificando la tarea,
las entradas y salidas, la arquitectura y los hiperparámetros. El objetivo
es que cada decisión de diseño sea justificable en términos de los
conceptos teóricos ya presentados.

%------------------------------------------------------------
\subsection{Definición del problema}
\label{subsec:def_problema}

El modelo implementado resuelve la tarea de \textit{modelado de lenguaje
a nivel carácter}: dada una secuencia de caracteres $c_1, c_2, \ldots,
c_t$, predecir la distribución de probabilidad sobre el siguiente carácter
$c_{t+1}$.

\begin{equation}
    P(c_{t+1} \mid c_1, c_2, \ldots, c_t)
\end{equation}

Esta tarea es representativa de los modelos de lenguaje modernos a gran
escala, reducida a una escala manejable. Operar a nivel carácter elimina
la necesidad de tokenizadores complejos: el vocabulario $\mathcal{V}$
está compuesto simplemente por todos los caracteres únicos del corpus de
entrenamiento (letras, espacios, signos de puntuación), con
$|\mathcal{V}|$ típicamente entre 60 y 100 símbolos.

La justificación de esta elección frente a alternativas más simples
(como clasificación) es que el modelado de lenguaje autoregresivo
requiere y ejercita todos los componentes del Transformer: embeddings,
codificación posicional, atención con máscara causal y proyección de
salida. Ningún componente es redundante.

%------------------------------------------------------------
\subsection{Representación de la entrada}
\label{subsec:entrada}

Cada carácter $c_t \in \mathcal{V}$ se representa como un índice entero
$v_t \in \{0, 1, \ldots, |\mathcal{V}|-1\}$ mediante un mapeo fijo
aprendido durante el preprocesamiento. Este índice se convierte en un
vector denso mediante la capa de embeddings:

\begin{equation}
    \bm{e}_t = E[v_t] \quad \in \mathbb{R}^{d_{\text{model}}}
\end{equation}

\noindent donde $E \in \mathbb{R}^{|\mathcal{V}| \times d_{\text{model}}}$
es la matriz de embeddings aprendida. A continuación se suma la
codificación posicional sinusoidal descrita en la
Sección~\ref{subsec:positional_encoding}:

\begin{equation}
    \bm{x}'_t = \bm{e}_t + PE_t
\end{equation}

La secuencia resultante $\bm{x}'_1, \ldots, \bm{x}'_T$ constituye la
entrada al encoder, donde $T$ es la longitud de contexto máxima.

%------------------------------------------------------------
\subsection{Arquitectura del modelo}
\label{subsec:arq_modelo}

El modelo sigue la variante \textit{decoder-only} con máscara causal,
que es la base arquitectónica de los modelos generativos modernos
(familia GPT). Se elige esta variante en lugar del encoder-decoder
original porque la tarea es de generación autoregresiva sobre una sola
secuencia, y un encoder-decoder requeriría dos secuencias distintas
(entrada y objetivo).

Formalmente, el modelo está compuesto por:

\begin{enumerate}
    \item \textbf{Capa de embedding:}
    $E \in \mathbb{R}^{|\mathcal{V}| \times d_{\text{model}}}$

    \item \textbf{Codificación posicional sinusoidal:}
    fija, no aprendida, sumada a los embeddings.

    \item \textbf{$N$ bloques encoder con máscara causal,}
    cada uno con:
    \begin{itemize}
        \item Multi-Head Self-Attention con máscara triangular inferior
        \item Feed-Forward Network con activación ReLU
        \item Conexiones residuales y Layer Normalization
    \end{itemize}

    \item \textbf{Proyección lineal de salida:}
    $W_o \in \mathbb{R}^{d_{\text{model}} \times |\mathcal{V}|}$,
    que mapea la representación final al espacio del vocabulario.

    \item \textbf{Softmax:} produce la distribución de probabilidad
    $\hat{P} \in \Delta^{|\mathcal{V}|-1}$ sobre el siguiente carácter.
\end{enumerate}

% ============================================================
% fig_pipeline_minimo.tex: Pipeline del Transformer mínimo
% Uso: % ============================================================
% fig_pipeline_minimo.tex: Pipeline del Transformer mínimo
% Uso: % ============================================================
% fig_pipeline_minimo.tex: Pipeline del Transformer mínimo
% Uso: \input{figures/fig_pipeline_minimo}
% ============================================================
\begin{figure}[H]
\centering
\begin{tikzpicture}[
    blk/.style={draw, rounded corners, minimum width=2.4cm,
                minimum height=0.7cm, font=\small, align=center},
    inp/.style={blk, fill=gray!10,    draw=gray!50},
    emb/.style={blk, fill=purple!10,  draw=purple!40},
    posblk/.style={blk, fill=teal!10,    draw=teal!50},
    enc/.style={blk, fill=blue!10,    draw=blue!40},
    outblk/.style={blk, fill=orange!10,  draw=orange!50},
    sft/.style={blk, fill=green!10,   draw=green!50!black},
    arrow/.style={-{Stealth}, thick},
    lbl/.style={font=\tiny, gray, align=center}
]

% ── Tokens de entrada ────────────────────────────────────────
\node[inp] (chars) at (0, 0) {``hola''\\{\tiny $[h,o,l,a]$}};
\node[lbl] at (0, -0.75) {caracteres};

% ── Índices enteros ──────────────────────────────────────────
\node[inp] (idx) at (2.8, 0) {índices\\{\tiny $[7,14,11,0]$}};
\node[lbl] at (2.8, -0.75) {vocabulario $|\mathcal{V}|$};

% ── Embedding ─────────────────────────────────────────────────
\node[emb] (emb) at (5.6, 0) {Embedding\\{\tiny $\mathbb{R}^{n \times d_{\text{model}}}$}};
\node[lbl] at (5.6, -0.75) {$E \in \mathbb{R}^{|\mathcal{V}| \times d}$};

% ── Positional Encoding ───────────────────────────────────────
\node[posblk] (pe) at (8.4, 0) {Positional\\Encoding};
\node[lbl] at (8.4, -0.75) {$\bm{x}'_t = \bm{e}_t + PE_t$};

% ── Encoder (N bloques) ───────────────────────────────────────
\node[enc] (enc) at (11.2, 0) {Encoder\\{\tiny $\times N$ bloques}};
\node[lbl] at (11.2, -0.75) {MHA + FFN + Norm};

% ── Proyección lineal ─────────────────────────────────────────
\node[outblk] (lin) at (11.2, -2.0) {Linear\\{\tiny $\mathbb{R}^{d} \to \mathbb{R}^{|\mathcal{V}|}$}};

% ── Softmax ──────────────────────────────────────────────────
\node[sft] (sft) at (8.4, -2.0) {Softmax};
\node[lbl] at (8.4, -2.75) {$P(\text{siguiente car.})$};

% ── Token predicho ────────────────────────────────────────────
\node[inp] (pred) at (5.6, -2.0) {``?''\\{\tiny token predicho}};
\node[lbl] at (5.6, -2.75) {argmax o muestreo};

% ── Flechas forward ───────────────────────────────────────────
\draw[arrow] (chars) -- (idx);
\draw[arrow] (idx)   -- (emb);
\draw[arrow] (emb)   -- (pe);
\draw[arrow] (pe)    -- (enc);
\draw[arrow] (enc)   -- (lin);
\draw[arrow] (lin)   -- (sft);
\draw[arrow] (sft)   -- (pred);

% ── Anotación: máscara causal ─────────────────────────────────
\node[draw=blue!20, fill=blue!4, rounded corners, font=\tiny,
      inner sep=4pt, align=left] at (11.2, 1.35)
    {M\'ascara causal:\\el token $t$ solo\\ve $1,\ldots,t{-}1$};
\draw[-{Stealth}, gray!40, thin] (11.2, 0.65) -- (11.2, 0.95);

% ── Objetivo de entrenamiento ─────────────────────────────────
\node[draw=orange!30, fill=orange!4, rounded corners, font=\tiny,
      inner sep=5pt, align=center] at (3.0, -2.0)
    {Entrenamiento:\\$\mathcal{L} = -\log P(c_{t+1} \mid c_1,\ldots,c_t)$\\
     Entrop\'ia cruzada sobre el\\siguiente car\'acter correcto};

% ── Flechas de flujo de la pérdida ────────────────────────────
\draw[-{Stealth}, orange!50, dashed, thick]
    (pred.west) -- +(-0.3, 0) -- +(-0.3, -0.7)
    node[below, font=\tiny, orange!70] {$\hat{y}$ vs $y$};

\end{tikzpicture}
\caption{Pipeline completo del Transformer m\'inimo implementado en este
trabajo. La secuencia de caracteres de entrada se convierte en \'indices
enteros, se proyecta al espacio de embeddings, se le suma la codificaci\'on
posicional, y se procesa a trav\'es de $N$ bloques encoder con m\'ascara
causal. La proyecci\'on lineal final y el softmax producen una distribuci\'on
de probabilidad sobre el vocabulario de caracteres. Durante el entrenamiento,
la p\'erdida de entrop\'ia cruzada compara la predicci\'on con el siguiente
car\'acter real de la secuencia.}
\label{fig:pipeline_minimo}
\end{figure}

% ============================================================
\begin{figure}[H]
\centering
\begin{tikzpicture}[
    blk/.style={draw, rounded corners, minimum width=2.4cm,
                minimum height=0.7cm, font=\small, align=center},
    inp/.style={blk, fill=gray!10,    draw=gray!50},
    emb/.style={blk, fill=purple!10,  draw=purple!40},
    posblk/.style={blk, fill=teal!10,    draw=teal!50},
    enc/.style={blk, fill=blue!10,    draw=blue!40},
    outblk/.style={blk, fill=orange!10,  draw=orange!50},
    sft/.style={blk, fill=green!10,   draw=green!50!black},
    arrow/.style={-{Stealth}, thick},
    lbl/.style={font=\tiny, gray, align=center}
]

% ── Tokens de entrada ────────────────────────────────────────
\node[inp] (chars) at (0, 0) {``hola''\\{\tiny $[h,o,l,a]$}};
\node[lbl] at (0, -0.75) {caracteres};

% ── Índices enteros ──────────────────────────────────────────
\node[inp] (idx) at (2.8, 0) {índices\\{\tiny $[7,14,11,0]$}};
\node[lbl] at (2.8, -0.75) {vocabulario $|\mathcal{V}|$};

% ── Embedding ─────────────────────────────────────────────────
\node[emb] (emb) at (5.6, 0) {Embedding\\{\tiny $\mathbb{R}^{n \times d_{\text{model}}}$}};
\node[lbl] at (5.6, -0.75) {$E \in \mathbb{R}^{|\mathcal{V}| \times d}$};

% ── Positional Encoding ───────────────────────────────────────
\node[posblk] (pe) at (8.4, 0) {Positional\\Encoding};
\node[lbl] at (8.4, -0.75) {$\bm{x}'_t = \bm{e}_t + PE_t$};

% ── Encoder (N bloques) ───────────────────────────────────────
\node[enc] (enc) at (11.2, 0) {Encoder\\{\tiny $\times N$ bloques}};
\node[lbl] at (11.2, -0.75) {MHA + FFN + Norm};

% ── Proyección lineal ─────────────────────────────────────────
\node[outblk] (lin) at (11.2, -2.0) {Linear\\{\tiny $\mathbb{R}^{d} \to \mathbb{R}^{|\mathcal{V}|}$}};

% ── Softmax ──────────────────────────────────────────────────
\node[sft] (sft) at (8.4, -2.0) {Softmax};
\node[lbl] at (8.4, -2.75) {$P(\text{siguiente car.})$};

% ── Token predicho ────────────────────────────────────────────
\node[inp] (pred) at (5.6, -2.0) {``?''\\{\tiny token predicho}};
\node[lbl] at (5.6, -2.75) {argmax o muestreo};

% ── Flechas forward ───────────────────────────────────────────
\draw[arrow] (chars) -- (idx);
\draw[arrow] (idx)   -- (emb);
\draw[arrow] (emb)   -- (pe);
\draw[arrow] (pe)    -- (enc);
\draw[arrow] (enc)   -- (lin);
\draw[arrow] (lin)   -- (sft);
\draw[arrow] (sft)   -- (pred);

% ── Anotación: máscara causal ─────────────────────────────────
\node[draw=blue!20, fill=blue!4, rounded corners, font=\tiny,
      inner sep=4pt, align=left] at (11.2, 1.35)
    {M\'ascara causal:\\el token $t$ solo\\ve $1,\ldots,t{-}1$};
\draw[-{Stealth}, gray!40, thin] (11.2, 0.65) -- (11.2, 0.95);

% ── Objetivo de entrenamiento ─────────────────────────────────
\node[draw=orange!30, fill=orange!4, rounded corners, font=\tiny,
      inner sep=5pt, align=center] at (3.0, -2.0)
    {Entrenamiento:\\$\mathcal{L} = -\log P(c_{t+1} \mid c_1,\ldots,c_t)$\\
     Entrop\'ia cruzada sobre el\\siguiente car\'acter correcto};

% ── Flechas de flujo de la pérdida ────────────────────────────
\draw[-{Stealth}, orange!50, dashed, thick]
    (pred.west) -- +(-0.3, 0) -- +(-0.3, -0.7)
    node[below, font=\tiny, orange!70] {$\hat{y}$ vs $y$};

\end{tikzpicture}
\caption{Pipeline completo del Transformer m\'inimo implementado en este
trabajo. La secuencia de caracteres de entrada se convierte en \'indices
enteros, se proyecta al espacio de embeddings, se le suma la codificaci\'on
posicional, y se procesa a trav\'es de $N$ bloques encoder con m\'ascara
causal. La proyecci\'on lineal final y el softmax producen una distribuci\'on
de probabilidad sobre el vocabulario de caracteres. Durante el entrenamiento,
la p\'erdida de entrop\'ia cruzada compara la predicci\'on con el siguiente
car\'acter real de la secuencia.}
\label{fig:pipeline_minimo}
\end{figure}

% ============================================================
\begin{figure}[H]
\centering
\begin{tikzpicture}[
    blk/.style={draw, rounded corners, minimum width=2.4cm,
                minimum height=0.7cm, font=\small, align=center},
    inp/.style={blk, fill=gray!10,    draw=gray!50},
    emb/.style={blk, fill=purple!10,  draw=purple!40},
    posblk/.style={blk, fill=teal!10,    draw=teal!50},
    enc/.style={blk, fill=blue!10,    draw=blue!40},
    outblk/.style={blk, fill=orange!10,  draw=orange!50},
    sft/.style={blk, fill=green!10,   draw=green!50!black},
    arrow/.style={-{Stealth}, thick},
    lbl/.style={font=\tiny, gray, align=center}
]

% ── Tokens de entrada ────────────────────────────────────────
\node[inp] (chars) at (0, 0) {``hola''\\{\tiny $[h,o,l,a]$}};
\node[lbl] at (0, -0.75) {caracteres};

% ── Índices enteros ──────────────────────────────────────────
\node[inp] (idx) at (2.8, 0) {índices\\{\tiny $[7,14,11,0]$}};
\node[lbl] at (2.8, -0.75) {vocabulario $|\mathcal{V}|$};

% ── Embedding ─────────────────────────────────────────────────
\node[emb] (emb) at (5.6, 0) {Embedding\\{\tiny $\mathbb{R}^{n \times d_{\text{model}}}$}};
\node[lbl] at (5.6, -0.75) {$E \in \mathbb{R}^{|\mathcal{V}| \times d}$};

% ── Positional Encoding ───────────────────────────────────────
\node[posblk] (pe) at (8.4, 0) {Positional\\Encoding};
\node[lbl] at (8.4, -0.75) {$\bm{x}'_t = \bm{e}_t + PE_t$};

% ── Encoder (N bloques) ───────────────────────────────────────
\node[enc] (enc) at (11.2, 0) {Encoder\\{\tiny $\times N$ bloques}};
\node[lbl] at (11.2, -0.75) {MHA + FFN + Norm};

% ── Proyección lineal ─────────────────────────────────────────
\node[outblk] (lin) at (11.2, -2.0) {Linear\\{\tiny $\mathbb{R}^{d} \to \mathbb{R}^{|\mathcal{V}|}$}};

% ── Softmax ──────────────────────────────────────────────────
\node[sft] (sft) at (8.4, -2.0) {Softmax};
\node[lbl] at (8.4, -2.75) {$P(\text{siguiente car.})$};

% ── Token predicho ────────────────────────────────────────────
\node[inp] (pred) at (5.6, -2.0) {``?''\\{\tiny token predicho}};
\node[lbl] at (5.6, -2.75) {argmax o muestreo};

% ── Flechas forward ───────────────────────────────────────────
\draw[arrow] (chars) -- (idx);
\draw[arrow] (idx)   -- (emb);
\draw[arrow] (emb)   -- (pe);
\draw[arrow] (pe)    -- (enc);
\draw[arrow] (enc)   -- (lin);
\draw[arrow] (lin)   -- (sft);
\draw[arrow] (sft)   -- (pred);

% ── Anotación: máscara causal ─────────────────────────────────
\node[draw=blue!20, fill=blue!4, rounded corners, font=\tiny,
      inner sep=4pt, align=left] at (11.2, 1.35)
    {M\'ascara causal:\\el token $t$ solo\\ve $1,\ldots,t{-}1$};
\draw[-{Stealth}, gray!40, thin] (11.2, 0.65) -- (11.2, 0.95);

% ── Objetivo de entrenamiento ─────────────────────────────────
\node[draw=orange!30, fill=orange!4, rounded corners, font=\tiny,
      inner sep=5pt, align=center] at (3.0, -2.0)
    {Entrenamiento:\\$\mathcal{L} = -\log P(c_{t+1} \mid c_1,\ldots,c_t)$\\
     Entrop\'ia cruzada sobre el\\siguiente car\'acter correcto};

% ── Flechas de flujo de la pérdida ────────────────────────────
\draw[-{Stealth}, orange!50, dashed, thick]
    (pred.west) -- +(-0.3, 0) -- +(-0.3, -0.7)
    node[below, font=\tiny, orange!70] {$\hat{y}$ vs $y$};

\end{tikzpicture}
\caption{Pipeline completo del Transformer m\'inimo implementado en este
trabajo. La secuencia de caracteres de entrada se convierte en \'indices
enteros, se proyecta al espacio de embeddings, se le suma la codificaci\'on
posicional, y se procesa a trav\'es de $N$ bloques encoder con m\'ascara
causal. La proyecci\'on lineal final y el softmax producen una distribuci\'on
de probabilidad sobre el vocabulario de caracteres. Durante el entrenamiento,
la p\'erdida de entrop\'ia cruzada compara la predicci\'on con el siguiente
car\'acter real de la secuencia.}
\label{fig:pipeline_minimo}
\end{figure}


%------------------------------------------------------------
\subsection{Función de pérdida y entrenamiento}
\label{subsec:entrenamiento}

El modelo se entrena minimizando la entropía cruzada entre la
distribución predicha $\hat{P}(c_{t+1})$ y el carácter real $c_{t+1}$:

\begin{equation}
    \mathcal{L} = -\frac{1}{T} \sum_{t=1}^{T}
    \log \hat{P}(c_{t+1} \mid c_1, \ldots, c_t)
\end{equation}

Esta pérdida se minimiza mediante descenso del gradiente con el
optimizador Adam~\cite{goodfellow2016deeplearning}, que extiende el
descenso del gradiente estocástico con estimaciones adaptativas de
momentos de primer y segundo orden, acelerando la convergencia en la
práctica.

Durante el entrenamiento, la máscara causal garantiza que el modelo no
acceda a información futura: el cómputo de $\hat{P}(c_{t+1})$ solo
utiliza $c_1, \ldots, c_t$. Esto hace posible calcular la pérdida para
todas las posiciones de la secuencia en paralelo, sin procesar los tokens
uno a uno como haría una RNN.

%------------------------------------------------------------
\subsection{Especificación formal de hiperparámetros}
\label{subsec:hiperparametros_minimo}

La Tabla~\ref{tab:hiperparametros_minimo} presenta la especificación completa
del modelo mínimo. Los valores están elegidos para que el entrenamiento
sea viable en hardware doméstico (CPU o GPU de consumo) en minutos, sin
sacrificar la capacidad de demostrar las propiedades fundamentales del
Transformer.

\begin{table}[H]
\centering
\caption{Especificación formal del Transformer mínimo implementado.}
\label{tab:hiperparametros_minimo}
\begin{tabularx}{\textwidth}{llXl}
\toprule
\textbf{Hiperparámetro} & \textbf{Notación} & \textbf{Descripción} & \textbf{Valor} \\
\midrule
Dimensión del modelo   & $d_{\text{model}}$ & Tamaño de embeddings y representaciones internas & 128 \\
Número de cabezas      & $h$               & Cabezas de atención paralelas & 4 \\
Dimensión por cabeza   & $d_k = d_v$       & $d_{\text{model}} / h$ & 32 \\
Número de bloques      & $N$               & Bloques encoder apilados & 2 \\
Dimensión FFN          & $d_{\text{ff}}$   & Capa interna de la red feed-forward & 512 \\
Longitud de contexto   & $T$               & Tokens procesados simultáneamente & 64 \\
Tasa de abandono       & $p_{\text{drop}}$ & Dropout para regularización & 0.1 \\
Tamaño de batch        & $B$               & Secuencias por paso de actualización & 32 \\
Tasa de aprendizaje    & $\eta$            & Paso del optimizador Adam & $3 \times 10^{-4}$ \\
\bottomrule
\end{tabularx}
\end{table}

Estos valores siguen el principio de que cada hiperparámetro debe ser
justificable. En particular:

\begin{itemize}
    \item $d_{\text{model}} = 128$ es suficiente para que el modelo
    aprenda representaciones de caracteres con estructura semántica,
    sin requerir memoria excesiva.

    \item $h = 4$ cabezas permiten capturar cuatro tipos de relaciones
    distintas en paralelo (por ejemplo, patrones fonéticos, morfológicos,
    sintácticos locales y dependencias de posición).

    \item $N = 2$ bloques producen un modelo con dos niveles de
    abstracción, lo que es suficiente para evidenciar el comportamiento
    jerárquico del Transformer.

    \item $T = 64$ garantiza un contexto suficiente para aprender
    patrones de longitud de palabra y estructura de oraciones simples.

    \item $d_{\text{ff}} = 512 = 4 \times d_{\text{model}}$ sigue la
    proporción estándar propuesta por Vaswani et al.~\cite{vaswani2017attention}.
\end{itemize}

%------------------------------------------------------------
\subsection{Comparación con modelos de producción}
\label{subsec:comparacion_escala}

Para contextualizar el modelo mínimo dentro del panorama actual, la
Tabla~\ref{tab:escala} lo compara con modelos de lenguaje representativos.

\begin{table}[H]
\centering
\caption{Comparación de escala entre el prototipo mínimo y modelos representativos.}
\label{tab:escala}
\begin{tabularx}{\textwidth}{lXXXX}
\toprule
\textbf{Modelo} & \textbf{$d_{\text{model}}$} & \textbf{$N$} & \textbf{$h$} & \textbf{Parámetros aprox.} \\
\midrule
Este trabajo (mínimo) & 128  & 2  & 4   & $\sim$0.41M \\
GPT-2 small           & 768  & 12 & 12  & $\sim$117M \\
GPT-2 large           & 1280 & 36 & 20  & $\sim$774M \\
GPT-3                 & 12288 & 96 & 96 & $\sim$175B \\
\bottomrule
\end{tabularx}
\end{table}

La diferencia de escala es enorme, pero los componentes son los mismos.
El prototipo pertenece a la misma familia decoder-only que GPT-2: atención
multi-cabeza con máscara causal, red feed-forward, conexiones residuales y
normalización de capa. Difiere en detalles de implementación (GPT-2 usa
posiciones aprendidas, activación GELU y normalización previa a cada
sub-capa), pero no en la idea. Esta es la propiedad más importante del
Transformer como arquitectura: su diseño escala de forma predecible y sin
cambios estructurales desde los $\sim$0.41M parámetros de este prototipo
(calculados en \texttt{verify/parametros.py} para un vocabulario de 65
caracteres) hasta los 175B de GPT-3~\cite{vaswani2017attention}.

En este sentido, el prototipo no es una simplificación del Transformer
real, sino el mismo modelo operando a escala didáctica.
